\chapter{M003 --- Surface Termination}

\section{Stage 1 --- Mathematical Specification}

\subsection{Scientific Purpose}

Construct a physically distinct surface termination from a primitive
surface.

\subsection{Transformation}

\[
\Phi_{\mathrm{termination}}:
\mathcal{S}_{\mathrm{primitive}}
\rightarrow
\mathcal{T}_{\mathrm{surface}}
\]

\subsection{Mathematical Inputs}

\begin{itemize}
\item Primitive surface
\item Bulk crystal motif
\end{itemize}

\subsection{Mathematical Outputs}

\begin{itemize}
\item Surface termination
\end{itemize}

\subsection{Domain}

Primitive surface space.

\subsection{Codomain}

Surface termination space.

\subsection{Preconditions}

\begin{itemize}
\item Valid primitive surface.
\item Well-defined surface orientation.
\end{itemize}

\subsection{Postconditions}

A surface termination representing a specific truncation of the bulk
crystal.

\subsection{Invariants}

\begin{itemize}
\item In-plane lattice.
\item Surface orientation.
\item Bulk crystallographic identity.
\end{itemize}

\subsection{Gauge Freedoms}

Lateral translations preserving the exposed termination.

\subsection{Non-Gauge Distinctions}

Different exposed atomic layers, stoichiometries, or polar terminations
are distinct scientific objects rather than alternative
representations.
